\scanpage{189}
Le calcul ``\unsure{bébête}'', en développant le déterminant et tenant compte des symétries (pour regrouper les termes qui se correspondent par \deleted{symétrie} op. de $\mathfrak{S}_3$) et $\Sigma^2$,
\[
\Sigma^2-\Delta^2=4\Pi(b,b',b'',t,t',t'',c,c',c'')
\]
avec
\[
\boxed{\begin{aligned}
\Pi={}&tt't''+(ct^2+c't'^2+c''t''^2)+(b^2c'c''+b'^2c''c+b''^2cc')\\
&-(b'b''ct+b''bc't'+bb'c''t'')
\end{aligned}}
\]
\inserted{\[
=tt't''+ct(t-b'b'')+c't'(t'-b''b)+c''t''(t''-bb')+(b^2c'c''+b'^2c''c+b''^2cc')
\]}
On aura donc
\[
\left\{\begin{aligned}
\tau+\tau'&=\Sigma\qquad(=bt+b't'+b''t''-bb'b'')\\
\tau\tau'&=\Pi\qquad\bigl(=tt't''+(ct^2+c't'^2+c''t''^2)\\
&\qquad\qquad\quad+(b^2c'c''+b'^2c''c+b''^2cc')\\
&\qquad\qquad\quad-(b'b''ct+b''bc't'+bb'c''t'')\bigr)
\end{aligned}\right.
\]
\[
\tau-\tau'=\varepsilon_0\Delta
\]
\margin{188}
